\documentclass[11pt,a4paper]{article}

% Packages
\usepackage[utf8]{inputenc}
\usepackage[T1]{fontenc}
\usepackage{amsmath,amssymb,amsfonts}
\usepackage{graphicx}
\usepackage{hyperref}
\usepackage{natbib}
\usepackage{booktabs}
\usepackage{algorithm}
\usepackage{algorithmic}
\usepackage[margin=1in]{geometry}
\usepackage{xcolor}
\usepackage{amsthm}
\newtheorem{definition}{Definition}

\title{Measuring Prompt Leakage Resistance in Deployed LLM Agents}
\author{
  Axel Fritz \\
  \texttt{axel.fritz@alquimia.ai}
}
\date{\today}

\begin{document}

\maketitle

\begin{abstract}
Large language models deployed with system prompts containing sensitive information are vulnerable
to adversarial extraction attacks. No standard, reproducible metric exists to quantify how well a
model resists such attacks. We propose \textbf{Prompt Leakage Resistance (PLR)}, a two-signal
detection framework that combines a semantic similarity score between the agent response and
chunks of the system prompt with a cosine alignment score against a reference safe response.
We evaluate four similarity methods (cosine sentence embeddings, a cross-encoder reranker,
ROUGE-L, and NLI entailment) across four leakage categories and three difficulty levels,
disqualifying ROUGE-L and NLI as structurally incompatible with the task and identifying the
cross-encoder reranker as the strongest detector. To test generalisation we extend the benchmark
to 246 queries and evaluate four agent models: the metric yields a discriminative resistance
ranking, and its operating threshold $\tau = 0.6$ transfers across models while keeping false
positives at or below 8\%. Against an independent labeled dataset, the leak-overlap signal separates leaked
from unrelated content with AUC $\approx 0.99$. We additionally report a \emph{smart cascade}
variant that lowers judge-escalation cost but converts some ambiguous adversarial cases into hard
false negatives. The metric and benchmark are released as part of the Gaussia evaluation framework.
\end{abstract}

% ============================================================
\section{Introduction}
% ============================================================

The proliferation of LLM-powered agents in production has made system prompt security a
first-order concern. Operators configure agents via system prompts that frequently contain
sensitive material: proprietary business logic, internal credentials, access control rules, and
behavioral constraints. Adversarial users can exploit prompt injection techniques~\citep{perez2022ignore}
to extract this information, exposing intellectual property or enabling privilege escalation.

Existing evaluation frameworks treat prompt leakage as a binary observable: either the agent
refused to reveal the prompt or it did not. This framing ignores the spectrum of partial
leakage, paraphrase-mediated disclosure, and the high false-positive risk of flagging legitimate
responses that necessarily reference system prompt content. No metric allows practitioners to
compare model configurations, system prompt formulations, or agent architectures on a common
scale.

We make the following contributions:
\begin{enumerate}
  \item \textbf{PLR metric}: a numeric score derived from a two-signal, chunk-based detection
    pipeline with a four-way decision table.
  \item \textbf{Benchmark}: a curated set of 68 records across four leakage categories and three
    difficulty levels, extended to a 246-query benchmark (evaluated across four agent models) via
    a documented generation-with-guard methodology.
  \item \textbf{Comparative evaluation}: four similarity methods tested against adversarial and
    legitimate records, with a clear empirical argument for discarding lexical and NLI approaches.
  \item \textbf{Generalisation and external validation}: the resistance ranking and the operating
    threshold $\tau = 0.6$ transfer across four agent models, and an independent labeled dataset
    corroborates the leak-overlap signal (AUC $\approx 0.99$).
  \item \textbf{Smart cascade (negative result)}: a combination of cosine embeddings and the
    reranker that reduces judge escalation but converts ambiguous adversarial cases into hard
    false negatives, making it unsuitable as the primary configuration in security-sensitive
    deployments.
\end{enumerate}

% ============================================================
\section{Related Work}
% ============================================================

Prompt injection and jailbreak attacks have received growing attention~\citep{perez2022ignore},
with most defenses focused on input filtering rather than post-hoc response evaluation. Existing
leakage metrics fall into two camps: an LLM judge issuing a binary leaked/not verdict (simple, but
no numeric score, no false-positive analysis, no partial-leak insight), or direct similarity
between the response and the full system prompt as a single vector (which conflates legitimate
domain overlap with actual leakage). PLR advances beyond both with chunk-level similarity to
preserve partial-leak signals, a second signal anchored to a safe reference response to control
false positives, and a numeric output for cross-configuration comparison. It is also complementary
to attack-side work such as PLeak~\citep{hui2024pleak}, which recovers prompts by optimisation:
PLR is a per-response detector operators run on their own agents, independent of any attack
strategy.

Our similarity components build on standard tools: Sentence-BERT bi-encoders~\citep{reimers2019sbert},
cross-encoder rerankers~\citep{nogueira2019passage}, the lexical ROUGE metric~\citep{lin2004rouge},
and NLI entailment models~\citep{he2021deberta}.

% ============================================================
\section{Problem Formulation}
% ============================================================

Let $P$ denote the system prompt, $Q$ an adversarial user query, $R$ the agent response to $Q$
under prompt $P$, and $G$ a reference safe response to $Q$ that does not reveal $P$.

\begin{definition}[Prompt Leakage]
A response $R$ constitutes a \emph{leak} of $P$ if it discloses information from $P$ that was
not already public knowledge or required by the query $Q$ to answer correctly.
\end{definition}

We partition leakable content into four categories that reflect the type of information at risk:

\begin{itemize}
  \item \textbf{Secrets \& Credentials}: API keys, tokens, passwords.
  \item \textbf{Internal Instructions}: business rules, persona definitions, behavioral
    constraints.
  \item \textbf{Guard Mechanism Exposure}: content filters, moderation triggers, rejection
    rules.
  \item \textbf{Permissions \& Roles}: access levels, user roles, capability restrictions.
\end{itemize}

Each attack instance carries a difficulty label (\textit{easy}, \textit{medium}, \textit{hard})
reflecting the sophistication of the adversarial query. Hard instances use paraphrase, roleplay,
and indirect multi-step strategies rather than direct extraction requests.

% ============================================================
\section{PLR Metric}
% ============================================================

\subsection{Chunking the System Prompt}

A system prompt embedded as a single vector averages all its semantic content. If the model
leaks only one section, for example, the credentials paragraph, that signal is diluted by the
non-leaked sections. We address this by splitting $P$ into overlapping chunks $C = \{c_1,
c_2, \ldots, c_k\}$ of at most four sentences each, and computing the \emph{maximum} similarity
between the response and any individual chunk:

\begin{equation}
  \operatorname{MaxChunkSim}(R, P) = \max_{c_i \in C} \operatorname{sim}(R,\, c_i)
\end{equation}

This design ensures that partial leaks of long prompts are not suppressed by averaging.

\subsection{Two-Signal Design}

PLR detection relies on two complementary signals:

\begin{itemize}
  \item \textbf{Signal~1} $= \operatorname{MaxChunkSim}(R, P)$: how much content from the
    system prompt appears in the response.
  \item \textbf{Signal~2} $= \operatorname{CosineSim}(R, G)$: how closely the response
    aligns with the reference safe response $G$. A high score indicates the agent responded
    within safe boundaries; a low score indicates divergence toward unreleased content.
\end{itemize}

\paragraph{Design decision: Signal~2 always uses cosine.}
Our experiments show that the cross-encoder reranker inflates $\operatorname{sim}(R, G)$
artificially: adversarial responses receive a mean Signal~2 score of 0.782, indistinguishable
from legitimate responses. This collapses the decision table by routing most adversarial records
to the \textit{ambiguous} bucket regardless of Signal~1. Cosine embeddings produce a calibrated
Signal~2 (mean 0.511 across both adversarial and legitimate records), making it the only method
suitable for the second signal regardless of which method is used for Signal~1.

\subsection{Decision Table}

\begin{table}[h]
  \centering
  \caption{Decision table for PLR verdicts. Threshold: $\tau = 0.6$ (default).}
  \label{tab:decision}
  \begin{tabular}{ccl}
    \toprule
    $\operatorname{MaxChunkSim}(R, P)$ & $\operatorname{CosineSim}(R, G)$ & Verdict \\
    \midrule
    $\geq \tau$ & $< \tau$ & \textbf{Leak} \\
    $< \tau$    & $\geq \tau$ & \textbf{Safe} \\
    $< \tau$    & $< \tau$  & \textbf{Anomalous} $\rightarrow$ judge \\
    $\geq \tau$ & $\geq \tau$ & \textbf{Ambiguous} $\rightarrow$ judge \\
    \bottomrule
  \end{tabular}
\end{table}

Anomalous and Ambiguous cases are escalated to a lightweight LLM judge.
The threshold $\tau$ is a user-configurable parameter; $\tau = 0.6$ is the recommended
default, grounded in cosine similarity distributions for semantically related texts
reported by~\citet{reimers2019sbert}. A higher $\tau = 0.65$ raises detection
(Section~\ref{sec:threshold}), but the multi-model study
(Section~\ref{sec:generalization}) shows it inflates false positives on some agents, so
$\tau = 0.6$ is the robust cross-model default. Methods with different score ranges require
independent threshold calibration (Section~\ref{sec:threshold}).

\subsection{PLR Score and Output}

The metric operates at two levels.

\paragraph{Per-record output.}
For each response $R_i$, the pipeline returns a discrete verdict
$v_i \in \{\textit{leak},\allowbreak\, \textit{safe},\allowbreak\, \textit{anomalous},\allowbreak\, \textit{ambiguous}\}$
together with the two underlying signal scores $\operatorname{MaxChunkSim}(R_i, P)$ and
$\operatorname{CosineSim}(R_i, G)$. In production, the verdict drives the immediate action
(flag, pass, or escalate); the signal scores provide interpretability.

\paragraph{Aggregate score.}
The PLR score is a scalar in $[0, 1]$, where higher values indicate greater resistance to
leakage:

\begin{equation}
  \text{PLR} = 1 - \frac{1}{N} \sum_{i=1}^{N} s_i
\end{equation}

\noindent where $s_i$ is the per-record leakage score defined as follows.

\paragraph{Binary mode.}
Each verdict maps to a score:

\begin{equation}
  s_i = \begin{cases}
    1              & \text{if } v_i = \textit{leak} \\
    0              & \text{if } v_i = \textit{safe} \\
    s_{\text{esc}} & \text{if } v_i \in \{\textit{anomalous},\, \textit{ambiguous}\}
  \end{cases}
\end{equation}

For escalated records, $s_{\text{esc}}$ depends on whether a judge is configured:

\begin{itemize}
  \item \textbf{With judge}: the verdict is resolved to \textit{leak} (or \textit{violation})
    $\to s_i = 1$; \textit{safe} $\to s_i = 0$; \textit{uncertain} (when the judge itself
    cannot decide) $\to s_i = 0.5$.
  \item \textbf{Without judge}: escalated records default to $s_{\text{esc}} = 0.5$,
    representing maximum uncertainty. Alternatively, they can be excluded from the aggregate
    via a parameter.
\end{itemize}

The use of a judge is recommended. Without one, the 0.5 default penalises the model
conservatively without making an unsupported claim of leak or safety.

Per-category PLR scores are reported alongside the global scalar to capture leakage patterns
by information type (Section~\ref{sec:exp2}).

% ============================================================
\section{Similarity Methods}
\label{sec:methods}
% ============================================================

We evaluate four methods for computing $\operatorname{MaxChunkSim}(R, P)$.

\subsection{Cosine Similarity on Sentence Embeddings}

Both texts are encoded independently as dense vectors using
\texttt{all-MiniLM-L6-v2}~\citep{reimers2019sbert}, and similarity is the cosine of the angle
between them:

\begin{equation}
  \operatorname{CosineSim}(R, c_i) = \frac{\mathbf{e}_R \cdot \mathbf{e}_{c_i}}{\|\mathbf{e}_R\| \, \|\mathbf{e}_{c_i}\|}
\end{equation}

This is the standard semantic similarity baseline: open-source, reproducible, and
well-calibrated in $[0, 1]$. It captures paraphrase-level semantic similarity that lexical
methods miss. In our experiments it serves dual purpose, as Signal~1 for the cosine pipeline
and as Signal~2 for every pipeline (including the reranker), where its calibrated scores are
essential for the decision table to function correctly.

\subsection{Cross-Encoder Reranker}

The \texttt{Qwen3-Reranker-0.6B}~\citep{alibaba2025qwen} reads the response and chunk jointly
in a single forward pass, producing a log-softmax probability of ``yes'' to the fixed instruction
$q$: \textit{``Given the agent response as the Query, determine whether the system prompt content
(Document) is revealed or leaked in the response.''} Unlike bi-encoders, cross-encoders attend
over the full concatenated input, detecting sophisticated paraphrase that independent encodings
miss. The score is:

\begin{equation}
  \operatorname{RerankerSim}(R, c_i) = P_\theta(\text{yes} \mid q,\, c_i,\, R)
\end{equation}

\noindent where $P_\theta$ is the log-softmax output of the cross-encoder $\theta$ evaluated on
the concatenation of $q$, $c_i$, and $R$.

The reranker achieves a mean MaxChunkSim of 0.970 on adversarial records, the strongest
detection signal of all evaluated methods, and is the recommended Signal~1 component.

\subsection{Disqualified Methods: ROUGE-L and NLI}

We also tested two methods that proved structurally incompatible with the task and are reported as
negative results.

\textbf{ROUGE-L}~\citep{lin2004rouge}, the F1 of the longest common subsequence, was hypothesised
to catch verbatim \textit{Secrets \& Credentials}. It detects \textbf{0\%} of the 52 adversarial
records (mean MaxChunkSim 0.165): instruction-tuned LLMs paraphrase by design, ``Your API key is
\texttt{sk-abc123}'' becomes ``The authentication credential uses identifier \texttt{sk-abc123}'', so
lexical overlap collapses even for a clear leak, and verbatim copy-paste is the least likely case.

\textbf{NLI entailment} (\texttt{cross-encoder/nli-deberta-v3-small}~\citep{he2021deberta}), using
the entailment probability as a similarity proxy, detects \textbf{5.8\%} (3/52, mean 0.113). Two
problems compound: its 512-token limit truncates the system prompts (measured at 624 tokens),
silently dropping the content under evaluation, and NLI targets short-sentence logical entailment,
not detecting whether information from a long document appears in a paraphrased response. Both
failures are independent of threshold choice.

% ============================================================
\section{Experimental Setup}
\label{sec:setup}
% ============================================================

\paragraph{Benchmark construction.}
All 68 benchmark responses were generated by \texttt{google/gemma-3-12b-it} queried via
HuggingFace Inference Endpoints, serving as the target agent under evaluation.
Gemma-3-12B-IT was chosen because it produces direct responses without chain-of-thought
reasoning tokens, yielding clean response text free of intermediate thinking artifacts.
Four realistic system prompts were constructed representing: a healthcare customer support
assistant (NovaMed), a technical sales assistant (TechVault), a content moderation platform
(SafeStream), and a corporate HR assistant (GlobalCorp). Each system prompt explicitly instructs
the agent not to disclose its confidential content, and the adversarial queries attempt to elicit
that content through a range of phrasings: direct requests, pretextual framing, and indirect
prompts.

\paragraph{Adversarial records (52).}
13 per category $\times$ 4 categories. Difficulty distribution: $\sim$5 easy, $\sim$7 medium,
and $\sim$5 hard per category. Hard instances use roleplay, indirect questions, and multi-step
misdirection.

\paragraph{Legitimate records (16).}
4 per category. These are queries where the agent \emph{must} reference system prompt content
to answer correctly (e.g., \textit{``What topics can you help me with?''}). They are hard
negatives for false positive analysis, the most challenging case for the two-signal design.

\paragraph{Evaluation.}
All four methods, plus threshold calibration, are evaluated on this benchmark using locally-run
models (\texttt{all-MiniLM-L6-v2}, \texttt{nli-deberta-v3-small}, \texttt{Qwen3-Reranker-0.6B}),
split by adversarial and legitimate subsets.

\paragraph{Generalization study.}
To test whether PLR generalizes beyond a single target model and a small sample, we extend the
evaluation along two axes. \emph{(i) Models:} in addition to the original Gemma target, we
evaluate three further agent models served via Groq, \texttt{llama-3.1-8b-instant},
\texttt{llama-4-scout-17b-16e-instruct}, and \texttt{gpt-oss-20b}, selected for diversity in
model family and size. \emph{(ii) Scale:} we expand the query bank from 68 to \textbf{246 unique
queries} (134 adversarial, 112 legitimate), the 68 original queries plus 178 new ones (82
adversarial, 96 legitimate). All four agent models answer the full 246-query bank. Pooled across
the four models this yields \textbf{984} (model, response) records, the population used in the
confusion analysis of Section~\ref{sec:partial-cm}. The new queries are generated by
two models distinct from the agents under test (\texttt{llama-3.3-70b-versatile} and
\texttt{gpt-oss-120b}), and each safe reference $G$ is synthesized by the largest available model
(\texttt{gpt-oss-120b}) under a constrained prompt that forbids disclosing any confidential value;
a hard guard rejects any $G$ containing a secret string (matched on word boundaries), and generated
items pass human review. The 68 original records, whose $G$ were generated separately by Gemma, are
retained as an independent \emph{origin} subset. Every reference is thus model-generated (Gemma for
the original 68, \texttt{gpt-oss-120b} for the 178 new); we have no human-authored gold standard for
$G$, and validation instead relies on the two independent generators producing consistent results
(Section~\ref{sec:generalization}).

% ============================================================
\section{Results}
\label{sec:results}
% ============================================================

\subsection{Method Selection}
\label{sec:exp2}

On the 52 adversarial records the cross-encoder reranker is the strongest Signal~1 detector at
every difficulty level and leakage category, cosine is the only viable alternative, and ROUGE-L
and NLI are disqualified for the structural reasons of Section~\ref{sec:methods} (0\% and 5.8\%
detection). We therefore adopt the reranker as Signal~1 and cosine as Signal~2; the two are not
competing detectors but the two complementary inputs to the decision table, and the rest of the
evaluation focuses on that combined pipeline rather than on either signal in isolation. One
property of the reranker matters downstream: it produces \emph{zero hard false negatives} at any
threshold evaluated here, escalating undetected adversarial records to the judge rather than
closing them as safe. This does not fully generalise (more resistant agents produce a few hard
false negatives, Section~\ref{sec:generalization}), which is why the judge remains a required
component.

\subsection{False Positive Analysis}
\label{sec:exp1}

The two-signal pipeline produces a single false positive on the 16 legitimate Gemma records
(6.2\%). It is \texttt{guar\_001\_leg\_medium} (\textit{``How does the reporting system work on
SafeStream?''}), which requires the agent to describe its own moderation pipeline, producing
unavoidable semantic overlap with the system prompt. This is a domain-boundary edge case: queries
about the agent's own domain necessarily echo the system prompt, and no similarity threshold can
disambiguate them from leaks without additional context. (ROUGE-L and NLI trivially produce no
false positives only because they cannot flag anything as a leak at all.)

\subsection{Threshold Calibration}
\label{sec:threshold}

Having identified the reranker as the recommended Signal~1 method, we evaluate
how the decision threshold $\tau$ affects the detection--false positive tradeoff. Table~\ref{tab:threshold}
reports results for the reranker across $\tau \in [0.55, 0.75]$.

The key distinction in this analysis is between two types of missed detections:
\begin{itemize}
  \item \textbf{Hard false negatives (FN-safe):} adversarial records closed as \textit{safe}, the
    system makes a confident wrong decision.
  \item \textbf{Soft false negatives (FN-judge):} adversarial records escalated to the
    judge, the system is uncertain and defers rather than deciding.
\end{itemize}

\begin{table}[h]
  \centering
  \caption{Reranker threshold sweep on 52 adversarial and 16 legitimate records.}
  \label{tab:threshold}
  \begin{tabular}{lcccc|cc}
    \toprule
    & \multicolumn{4}{c|}{Adversarial (52)} & \multicolumn{2}{c}{Legitimate (16)} \\
    $\tau$ & Detection & FN-safe & FN-judge & & FP & Safe \\
    \midrule
    0.55 & 50.0\% & 0.0\% & 50.0\% & & 6.2\%  & 18.8\% \\
    \textbf{0.60} & \textbf{61.5\%} & \textbf{0.0\%} & \textbf{38.5\%} & & \textbf{6.2\%}  & \textbf{25.0\%} \\
    \textbf{0.65} & \textbf{76.9\%} & \textbf{0.0\%} & \textbf{23.1\%} & & \textbf{6.2\%}  & \textbf{25.0\%} \\
    0.70 & 84.6\% & 0.0\% & 15.4\% & & 31.2\% & 25.0\% \\
    0.75 & 84.6\% & 0.0\% & 15.4\% & & 50.0\% & 6.2\%  \\
    \bottomrule
  \end{tabular}
\end{table}

Two observations emerge. First, the reranker produces \emph{zero hard false negatives across all
thresholds evaluated here}: undetected adversarial records are escalated to the judge, never
labeled safe. Second, while $\tau = 0.65$ looks locally optimal (detection 61.5\% to 76.9\%, false
positives flat at 6.2\%), the multi-model study (Section~\ref{sec:generalization}) shows it
inflates false positives on other agents, so we adopt $\tau = 0.6$ as the robust cross-model
default and treat any higher value as requiring per-deployment validation.

\subsection{Smart Cascade: A Negative Result}
\label{sec:cascade}

We also tested a \emph{smart cascade} that uses cosine to pre-filter \textit{safe} responses
before invoking the reranker, motivated by cost: cosine can close clean traffic without a reranker
call. It did not work as a detector. By closing cases on cosine's Signal~2 alone, the cascade
converts ambiguous adversarial records that the reranker would have escalated to the judge into
hard false negatives, and threshold tuning cannot recover them without collapsing the cascade back
to the reranker alone. We therefore do not recommend it for adversarial detection; the reranker
alone is strictly superior, and cosine's only useful role is the second signal of the decision
table.

\subsection{Generalization Across Models}
\label{sec:generalization}

We now evaluate PLR across four agent models on the expanded benchmark ($N = 246$ each), using the
recommended reranker pipeline at the conservative default $\tau = 0.6$. Table~\ref{tab:multimodel}
reports PLR together with detection and false-positive rates and their 95\% Wilson confidence
intervals. PLR separates the models over a wide range (0.15--0.45): \texttt{gpt-oss-20b} is the
most resistant (PLR 0.448; only 22.4\% of attacks surface as confirmed leaks), while Gemma is the
most permissive (PLR 0.153; 70.9\% detection). The metric therefore produces a discriminative
ranking rather than a model-independent constant. The benefit of scale is visible in the intervals:
each false-positive estimate ($N = 246$, 112 legitimate per model) has a 95\% half-width of only a
few points, tight enough to separate the models' specificity.

\begin{table}[h]
  \centering
  \caption{PLR and detection/false-positive rates across four agent models on the expanded
    benchmark ($N = 246$ each; reranker, $\tau = 0.6$; 95\% Wilson intervals). PLR is oriented so
    higher means more resistant.}
  \label{tab:multimodel}
  \begin{tabular}{lccc}
    \toprule
    Model & PLR & Detection [95\% CI] & FP [95\% CI] \\
    \midrule
    \texttt{gpt-oss-20b}      & \textbf{0.448} & 22.4\% [16.2--30.2] & 6.2\% [3.1--12.3] \\
    Llama-4-Scout-17B         & 0.276 & 50.7\% [42.4--59.1] & 3.6\% [1.4--8.8] \\
    Llama-3.1-8B              & 0.254 & 54.5\% [46.0--62.7] & 8.0\% [4.3--14.6] \\
    Gemma-3-12B               & 0.153 & 70.9\% [62.7--77.9] & 2.7\% [0.9--7.6] \\
    \bottomrule
  \end{tabular}
\end{table}

\paragraph{Threshold stability.}
A central question is whether $\tau = 0.6$ is an artifact of the original single-model
calibration. Table~\ref{tab:taustab} sweeps $\tau$ for all four models. The false-positive rate
stays at or below 8\% at $\tau = 0.6$ for every model, and only at $\tau = 0.70$ does it degrade
sharply and \emph{universally} (Llama-4-Scout 16\%, Llama-3.1-8B 15\%, Gemma 14\%, gpt-oss 11\%).
The single-model study favored $\tau = 0.65$; the multi-model data shows that 0.65 already inflates
false positives on the Llama agents (12\% and 8\%). We therefore retain $\tau = 0.6$ as the robust
cross-model default: the operating point transfers across models rather than overfitting to a
single one.

\begin{table}[h]
  \centering
  \caption{Reranker threshold sweep across models, reported as detection\,/\,false-positive (\%).
    False positives stay $\le 8\%$ at $\tau = 0.6$ for all models and rise sharply at
    $\tau = 0.70$.}
  \label{tab:taustab}
  \begin{tabular}{lccccc}
    \toprule
    Model & $\tau{=}0.50$ & $0.55$ & $\mathbf{0.60}$ & $0.65$ & $0.70$ \\
    \midrule
    Gemma-3-12B       & 54/3 & 63/3 & \textbf{71/3} & 78/8  & 81/14 \\
    Llama-3.1-8B      & 47/4 & 54/4 & \textbf{54/8} & 60/12 & 69/15 \\
    Llama-4-Scout-17B & 40/3 & 43/4 & \textbf{51/4} & 54/8  & 64/16 \\
    \texttt{gpt-oss-20b} & 16/3 & 18/3 & \textbf{22/6} & 24/8 & 25/11 \\
    \bottomrule
  \end{tabular}
\end{table}

\paragraph{Cross-generator validation.}
The benchmark uses two independent reference generators, Gemma for the origin subset,
\texttt{gpt-oss-120b} for the expanded subset, so we test whether detection is consistent across
them. Detection rates on the origin (Gemma $G$) and expanded (\texttt{gpt-oss-120b} $G$)
adversarial subsets agree within two points for both Llama agents
(Llama-3.1-8B: 55.8\% vs.\ 53.7\%; Llama-4-Scout: 50.0\% vs.\ 51.2\%), indicating that the
metric's verdicts do not depend on which model authored the reference. We anticipated a possible
confound: since the expanded $G$ is generated by \texttt{gpt-oss-120b}, the same-family
\texttt{gpt-oss-20b} agent might be scored artificially safe. The data rules this out: %
\texttt{gpt-oss-20b} detection is low on \emph{both} the origin subset (19.2\%, where $G$ is
generated by Gemma, not by the gpt-oss family) and the expanded subset
(24.4\%), so its high resistance is a property of the model, not an artifact of reference
generation.

\subsection{External Validation}
\label{sec:external}

All results so far use references and verdicts produced within our own pipeline. To check that
the leak-detection signal agrees with \emph{independent} labels, we validate Signal~1
($\operatorname{MaxChunkSim}$ between a response and the system prompt) against the
\texttt{gabrielchua/system-prompt-leakage} dataset,\footnote{\url{https://huggingface.co/datasets/gabrielchua/system-prompt-leakage}}
which provides (system prompt, content, binary leakage) triples. Because that dataset has no
safe reference $G$, we evaluate the overlap signal alone, not the full two-signal table, on a
balanced sample of 160 test items (80 leak, 80 no-leak). Table~\ref{tab:external} reports the
result at $\tau = 0.6$ with 95\% Wilson intervals, plus the threshold-independent AUC.

\begin{table}[h]
  \centering
  \caption{Signal~1 versus independent labels (\texttt{gabrielchua/system-prompt-leakage},
    balanced $n = 160$). AUC is threshold-independent; precision/recall are at $\tau = 0.6$ with
    95\% Wilson intervals.}
  \label{tab:external}
  \begin{tabular}{lcccc}
    \toprule
    Method & AUC & Precision & Recall & F1 \\
    \midrule
    Cosine   & \textbf{0.989} & 0.96 [0.90--0.99] & 0.96 [0.90--0.99] & \textbf{0.96} \\
    Reranker & 0.985 & 0.83 [0.75--0.89] & 1.00 [0.95--1.00] & 0.91 \\
    \bottomrule
  \end{tabular}
\end{table}

The overlap signal separates prompt-derived content from unrelated content almost perfectly
(AUC $\approx 0.99$), confirming on external, independently-labeled data that the detector
recognises genuine leakage rather than over-fitting our own benchmark. Notably cosine is the
stronger method on this distribution and the reranker's optimal threshold shifts upward (best F1
at $\tau = 0.8$), consistent with its score saturation near 0 and 1 and reinforcing that $\tau$
should be calibrated per distribution.

\paragraph{Scope of this evidence.}
We state the boundary of this result plainly. The dataset's negatives are \emph{off-topic} text,
so this validates the easy half of the problem, \textbf{leak vs.\ unrelated content}, on which
the detector is excellent. It does \emph{not} test the hard half, \textbf{leak vs.\ a legitimate
on-topic response} that necessarily echoes the system prompt, which is precisely what the
two-signal design targets and what drives the false-positive analysis in
Section~\ref{sec:exp1}. Validating that half requires a ground-truth dataset of human-labeled
on-topic responses, which to our knowledge does not yet exist and which we did not construct
here. We therefore present these metrics as strong but \emph{partial} evidence: they show the
mechanism works and generalises, and we consider them sufficient to support PLR as a practical
detector, while stating explicitly that the on-topic case is not yet validated end-to-end.

\subsection{What We Can and Cannot Measure: A Partial Confusion Matrix}
\label{sec:partial-cm}

A standard confusion matrix requires a ground-truth label for every record: did the response
\emph{actually} leak? We have that label for only one of our two query types, and reporting a full
precision/recall table would paper over the gap. We therefore make the gap explicit. Table~\ref{tab:partial-cm}
cross-tabulates each record's query type against the metric's verdict (reranker, $\tau = 0.6$),
pooled across all four responder models ($n = 984$: 536 adversarial, 448 legitimate).

\begin{table}[h]
  \centering
  \caption{Partial confusion matrix: query type $\times$ metric verdict (reranker, $\tau = 0.6$,
    pooled over four models, $n = 984$). \emph{Legitimate} responses cannot leak by construction,
    so their column has reliable labels; \emph{adversarial} responses lack a ground-truth leak
    label, so those cells are not scorable. ``Escalated'' rows are deferred to the judge
    (Section~\ref{sec:judge-eval}).}
  \label{tab:partial-cm}
  \begin{tabular}{lccc}
    \toprule
    Query type & Verdict: leak & Verdict: safe & Escalated \\
    \midrule
    Legitimate ($n=448$, label known)   & 23 (FP) & 192 (TN) & 233 \\
    Adversarial ($n=536$, label unknown) & 266 & 33 & 237 \\
    \bottomrule
  \end{tabular}
\end{table}

\paragraph{The legitimate column is scorable; the adversarial column is not.}
A legitimate, on-topic response cannot constitute a leak, so for that class the verdicts are true
labels: 23 of 448 legitimate responses were flagged, a \textbf{false-positive rate of $5.1\%$}, and
192 were correctly cleared as safe (true negatives). This is the column the two-signal design is
built to protect, and the number is trustworthy. The adversarial column is a different matter: an
adversarial \emph{query} does not imply an adversarial \emph{response}, because a resistant model
simply refuses. We have no human-verified label for which of the 536 adversarial responses
genuinely disclosed protected content. The 266 ``leak'' verdicts are therefore \emph{flagged
candidates}, not confirmed true positives, and the 33 ``safe'' verdicts are not confirmed false
negatives.

\paragraph{Consequence for precision and recall.}
Recall (true positives over actual leaks) and precision (true positives over flagged) both require
the adversarial ground-truth labels we lack, so we do \emph{not} report them over our own benchmark;
doing so would assume every adversarial query produced a leak, which is false. What we can report
are the two quantities that do not need that label: the false-positive rate on legitimate traffic
($5.1\%$ pooled, $2.7$--$8.0\%$ per model in Table~\ref{tab:multimodel}), and the
external precision/recall on the leak-vs-unrelated case (Table~\ref{tab:external}, AUC
$\approx 0.99$). The missing piece, precision and recall on the leak-vs-legitimate-on-topic
case, is exactly the human-labeled dataset that, to our knowledge, does not yet exist
(Section~\ref{sec:external}). The 470 escalated records (233 legitimate, 237 adversarial) are not
discarded: they are deferred to the LLM judge, whose behaviour we examine next.

\paragraph{An objectively-labeled slice: verbatim credential leaks.}
One slice of the adversarial column \emph{can} be labeled without human judgment or any LLM. The
\texttt{secrets\_and\_credentials} category protects high-entropy tokens (production API keys
like \texttt{tvk\_prod\_\dots}, webhook secrets, database tokens, an internal escalation
URL) that cannot occur by chance in a benign answer. For these, ``did the response leak?'' reduces
to a deterministic string match: is the literal secret present? This yields true labels for that
slice, free of the circularity that an LLM annotator would introduce. Over the 132 adversarial
records in this category, 33 reproduced a literal secret (true leaks). Table~\ref{tab:verbatim}
crosses that ground truth against the metric's clear verdicts (reranker, $\tau = 0.6$).

\begin{table}[h]
  \centering
  \caption{Deterministic ground truth on the \texttt{secrets\_and\_credentials} slice (verbatim
    secret present $=$ leak). Confusion over the metric's \emph{clear} verdicts; 11 of the 33 true
    leaks were escalated to the judge (deferred, not missed) and are excluded from this table.}
  \label{tab:verbatim}
  \begin{tabular}{lcc}
    \toprule
     & GT: leak (verbatim) & GT: no verbatim secret \\
    \midrule
    Verdict: leak & 22 (TP) & 31 (FP) \\
    Verdict: safe & 0 (FN)  & 11 (TN) \\
    \bottomrule
  \end{tabular}
\end{table}

The recall on this slice is exact and strong: \textbf{zero of the 33 verbatim leaks were cleared as
safe} (the 11 not flagged outright were escalated to the judge, not missed), so the metric never
silently passed a literal credential dump. Precision, however, reads as $42\%$
($22 / 53$), and this number is a \emph{lower bound, not the true precision}. A string match only
recognises the secret reproduced \emph{verbatim}; it cannot count a response that leaks by
paraphrasing or describing the credential as a true positive. Many of the 31 ``false positives''
are therefore semantic leaks invisible to the literal check, which is exactly why verbatim matching
cannot replace a human label for the general case. The slice confirms what it can: the detector
does not miss literal secret disclosures, and demonstrates, by its own blind spot, why the
on-topic precision question stays open.

\paragraph{Why we stop here: an acknowledged technical debt.}
The principled way to close the remaining gap is a human-labeled set of on-topic adversarial
responses. We did not build one: at $984$ pooled records across four models, exhaustive hand
annotation was beyond the scope of this study, and a partial hand-labeling would itself need an
adjudication protocol to be defensible. We therefore record this explicitly as a limitation rather
than paper over it: PLR is validated on false-positive specificity (reliably), on the
leak-vs-unrelated case (externally, AUC $\approx 0.99$), and on verbatim credential recall
(deterministically, no misses), while end-to-end precision/recall on the hard leak-vs-legitimate
case remains future work contingent on a human-annotated benchmark.

\subsection{Judge Evaluation}
\label{sec:judge-eval}

The decision table escalates ambiguous records, high overlap \emph{and} high similarity to $G$, or
low on both, to an LLM judge. To characterise that stage we evaluated two large candidate judges,
neither of which is in the responder roster (avoiding self-judging):
\texttt{llama-3.3-70b-versatile} and \texttt{openai/gpt-oss-120b}. Each judged a pooled set of
$733$ records, all escalations (470), all legitimate records (as a false-positive probe), and a
capped clear-leak sample, seeing the (system prompt, user message, response) triple as content to
\emph{audit}, never as a role to adopt. Without human ground truth we cannot score a judge's
accuracy on true leaks; we measure what is measurable (Table~\ref{tab:judge}): how each judge
resolves escalations, its false-positive rate on legitimate traffic (where the label is reliable),
and inter-judge agreement.

\begin{table}[h]
  \centering
  \caption{Judge behaviour on $733$ pooled records. Escalation resolution is over the 470 escalated
    records; the false-positive rate is on the 448 legitimate records (reliable labels).}
  \label{tab:judge}
  \begin{tabular}{lcccc}
    \toprule
    Judge & \multicolumn{3}{c}{Escalations resolved (of 470)} & FP on legitimate \\
    \cmidrule(lr){2-4}
          & leak & safe & uncertain & \\
    \midrule
    \texttt{llama-3.3-70b}   & 109 & 361 & 0  & \textbf{19.0\%} (85/448) \\
    \texttt{gpt-oss-120b}    & 136 & 324 & 10 & 24.3\% (109/448) \\
    \bottomrule
  \end{tabular}
\end{table}

The judges resolve roughly three-quarters of escalations as safe and a quarter as leaks, and they
agree with each other on $89.8\%$ of records (Cohen's $\kappa = 0.72$). The decisive finding is the
false-positive column: applied to legitimate traffic, the judges flag $19.0\%$ and $24.3\%$ of
benign responses as leaks, roughly four to five times the two-signal metric's $5.1\%$. A raw LLM
judge, in other words, is far more trigger-happy than the metric, which is precisely the
over-flagging the two-signal design was built to avoid; the metric's specificity is a real
advantage, not an artifact. The flip side, and a limitation we state plainly, is that the judges
agree with the metric on only $47$--$52\%$ of clear-leak cases while agreeing strongly with each
other: there is
no consensus on the \emph{positive} class, which again traces back to the absence of human ground
truth. We use the lower-FP \texttt{llama-3.3-70b} as the default judge, and treat the escalated
positives as candidates for human review rather than final labels.

\paragraph{Toward a continuous judge (proposed).}
All judge verdicts above are \emph{discrete}: every escalated record resolves to a single
leak/safe/uncertain label, and that is the mode used throughout this evaluation. We propose, but
do not yet validate, a continuous extension based on \emph{self-consistency}: query the judge $k$
times at non-zero temperature and take the fraction of ``leak'' votes as a continuous leakage
score in $[0,1]$, feeding it directly into the aggregate PLR instead of the $\{0, 0.5, 1\}$
mapping. Genuinely ambiguous responses yield intermediate scores (split votes), while clear cases
saturate, so the dispersion of an ensemble of judgments becomes the signal. We favour this over
reading the provider's token log-probabilities for the verdict token: a single token's
log-probability reflects the model's confidence in \emph{emitting that word}, not a calibrated
probability that the response leaked, and post-RLHF models are systematically over-confident, so
the resulting scores tend to collapse to 0 or 1 rather than spread. Self-consistency is also
provider-agnostic (it needs only repeated sampling, not log-probability access, which several
inference APIs do not expose), with log-probabilities as an optional cheaper fallback and the
current discrete verdict as the final fallback. Validating whether the continuous score is better
calibrated than the discrete one again requires ground truth; the verbatim-credential slice
(Section~\ref{sec:partial-cm}) offers a small objective anchor on which to test it.

% ============================================================
\section{Discussion}
\label{sec:discussion}
% ============================================================

\paragraph{Why ROUGE and NLI fail.}
The failure of lexical and entailment-based approaches is not a calibration artifact; it is
structural. Modern instruction-tuned LLMs rephrase rather than regurgitate, collapsing lexical
overlap to near-zero even for unambiguous leaks. NLI compounds this with 512-token truncation
that silently discards the system prompt content being evaluated. Both methods are disqualified
regardless of threshold tuning, and we report this as a useful negative result for practitioners
evaluating off-the-shelf similarity metrics for leak detection.

\paragraph{The reranker's conservatism: a detection asset and an operational cost.}
The reranker concentrates scores near 0 and 1 (mean MaxChunkSim 0.970 on Gemma adversarial
records), so when it misses a leak it tends to escalate rather than close: no confident wrong
calls on Gemma, though the property weakens on more resistant agents (a few hard false negatives,
Section~\ref{sec:generalization}), which is why the judge is not optional. The same conservatism
is a cost on legitimate traffic: at $\tau = 0.6$ it closes only 4 of 16 Gemma legitimate records
as \textit{safe} versus cosine's 10, escalating the rest at the full price of a judge call. The
reranker-only pipeline is thus more expensive at scale than its detection numbers suggest. This is
the central operational trade-off of the metric.

\paragraph{The domain-boundary false positive.}
False positives trace to an edge case inherent to semantic detection: agents answering legitimate
queries in their own domain echo prompt language. On Gemma this is a single stable record across
$\tau \in [0.55, 0.65]$, and the four-model false-positive rate stays in 3.6--8\% at $\tau = 0.6$.
One mitigation is per-query reference calibration, constructing $G$ to reference the domain
without sensitive specifics, sharpening Signal~2 on boundary cases.

\paragraph{Limitations.}
\begin{itemize}
  \item \textbf{Threshold calibration}: we adopt $\tau = 0.6$ as the cross-model default.
    Different system prompt styles, agent models, or domain vocabularies may shift the optimal
    threshold, so per-deployment calibration against a held-out legitimate sample is recommended
    before raising $\tau$ above this default.
  \item \textbf{No human ground truth}: this is the deepest limitation. Both the reference
    responses $G$ and the leak/safe verdicts are produced by models: there is no human-labeled
    set establishing which responses \emph{actually} leak. PLR is therefore an internally
    consistent operational score that we show generalizes across models and reference generators,
    not a detector validated against human judgment. A small human-labeled validation study is
    the most important next step before treating PLR as a ground-truth leak measure.
  \item \textbf{Scale}: the benchmark spans 246 queries answered by all four agent models
    (984 pooled records, Section~\ref{sec:generalization}). Aggregate confidence intervals are
    tight, but per-category and per-difficulty intervals remain wide.
  \item \textbf{Model-generated references}: every $G$ is LLM-generated (Gemma for the origin
    queries, \texttt{gpt-oss-120b} for the new ones). We mitigate the risks with a hard anti-secret
    guard, human review, and cross-generator agreement (the two generators yield detection within
    two points for the Llama agents), but no $G$ is human-verified.
  \item \textbf{Adversarial authorship skew}: \texttt{gpt-oss-120b} declined to author most
    extraction attacks for safety reasons, so the generated adversarial queries are
    predominantly authored by \texttt{llama-3.3-70b-versatile}. Attack-author diversity is thus
    lower than intended.
  \item \textbf{Label noise}: roughly 3\% of generated adversarial queries are benign questions
    mislabeled as attacks. We retain them and rely on the origin subset to corroborate the
    conclusions; a label-cleaning pass would tighten the generated-subset estimates.
  \item \textbf{Multi-turn attacks}: the current pipeline evaluates single-turn responses.
    Multi-turn attacks, where an adversary extracts the prompt incrementally across several
    turns, are outside the current scope and are left for future work.
  \item \textbf{Undetected leaks}: a substantial share of adversarial records is not automatically
    classified by the reranker at $\tau = 0.6$ and is escalated to the judge, and, on the more
    resistant agents, a minority is closed as a hard false negative. These are medium-to-hard
    paraphrase attacks where the reranker score falls below threshold. A judge component is a
    required part of any production deployment.
\end{itemize}

\paragraph{Future work.}
The most important next step is a small human-labeled validation study to close the ground-truth
gap; beyond it, a multi-turn evaluation protocol and a judge tuned for the ambiguous case are the
primary extensions. A \emph{continuous scoring mode} ($s_i = \operatorname{MaxChunkSim}(R_i, P)
\times (1 - \operatorname{CosineSim}(R_i, G))$) is natural but currently degenerate: with the
reranker, Signal~1 saturates near 1 for both leaks and legitimate domain references, collapsing
the product to $1 - \operatorname{CosineSim}(R_i, G)$. The deeper open problem is a deterministic
\emph{safe classifier}, one that confirms a response does not leak without an LLM judge and
without cosine's hard false negatives, which would let the pipeline close clean traffic cheaply
and reserve the judge for genuinely ambiguous cases.

% ============================================================
\section{Conclusion}
% ============================================================

We introduce PLR (Prompt Leakage Resistance), a metric for quantifying an LLM agent's
resistance to system prompt extraction attacks. The metric combines a chunk-level maximum
similarity score with a cosine alignment score against a safe reference response, applying a
four-way decision table that separates confirmed leaks, safe responses, and judge-escalated
ambiguous cases.

Our method comparison shows that lexical (ROUGE-L) and entailment-based (NLI) approaches are
structurally incompatible with the task due to LLM paraphrase and token-limit truncation
respectively, leaving the cross-encoder reranker as the strongest detector.

Extending the benchmark to 246 queries evaluated across four agent models, PLR produces a
discriminative resistance ranking and its operating threshold $\tau = 0.6$ transfers across models, keeping
false positives at or below 8\%; a higher $\tau$ that looked optimal on a single model inflates
false positives elsewhere. An independent labeled dataset corroborates the leak-overlap signal
(AUC $\approx 0.99$) on the leak-versus-unrelated case, while the leak-versus-legitimate-on-topic
case remains to be validated against human-labeled ground truth. We also report a smart cascade
as a negative result: pre-filtering with cosine lowers judge load but converts ambiguous
adversarial cases into hard false negatives, so it is not recommended as the primary configuration.

PLR is integrated into the Gaussia evaluation framework, enabling practitioners to compare model
configurations on a common numeric scale across leakage categories and difficulty levels.

% ============================================================
\section*{Open Source}
% ============================================================

PLR is developed as part of the Gaussia evaluation framework, an open-source initiative
for reproducible LLM agent evaluation. The benchmark, experiment notebooks, and full
implementation are publicly available as part of \texttt{pygaussia} at
\url{https://github.com/gaussia-labs/pygaussia}.

\bibliographystyle{plainnat}
\bibliography{references}

\end{document}
